\subsection{Separately continuous bilinear mappings}

Lemma 1. --- Let E and F be two locally convex metrizable spaces, and u be a continuous linear mapping from \(E_{b}^{\prime}\) into F. Then there exists a neighbourhood U of 0 in \(E_{b}^{\prime}\) whose image under u is bounded in F.

Let \((\mathrm{U}_n)_{n\in \mathbf{N}}\) (resp. \((\mathrm{V}_n)_{n\in \mathbf{N}}\) ) be a fundamental system of neighbourhoods of 0 in E (resp. F). We assume that the sets \(\mathrm{U}_n\) are balanced and form a decreasing sequence. Since \(u\) is continuous, for every \(n\in \mathbf{N}\) , there exists a bounded set \(\mathrm{B}_n\) in E such that \(u(\mathrm{B}_n^\circ)\subset \mathrm{V}_n\) . Since \(\mathrm{B}_n\) is bounded, there exists a real number \(\lambda_n > 0\) such that \(\lambda_{n}\mathrm{B}_{n}\subset \mathrm{U}_{n}\) . Put \(\mathrm{B} = \bigcup_{n\in \mathbf{N}}\lambda_n\mathrm{B}_n\) .

We shall prove that the set B is bounded in E, in other words that for every integer \(m \geqslant 0\) , there exists a real number \(\mu > 0\) such that \(\mu B \subset U_{m}\) . Since the sets \(B_{n}\) are bounded, there exists a real number \(\mu\) such that \(0 < \mu \leqslant 1\) and such that \(\mu \cdot (\lambda_{n} B_{n}) \subset U_{m}\) for \(0 \leqslant n \leqslant m\) ; we have also \(\lambda_{n} B_{n} \subset U_{n} \subset U_{m}\) if n \textgreater{} m; hence \(\mu B \subset U_{m}\) since \(U_{m}\) is balanced.

Let \(\mathbf{U}\) be the polar of \(\mathbf{B}\) in \(\mathrm{E}_b'\) . This is a neighbourhood of 0 in \(\mathrm{E}_b'\) and we have \(\lambda_n\mathrm{B}^\circ \subset \mathrm{B}_n^\circ\) , hence \(\lambda_n u(\mathrm{U}) \subset \mathrm{V}_n\) for all \(n \in \mathbf{N}\) . Consequently \(u(\mathrm{U})\) is bounded in \(\mathrm{F}\) .

THEOREM 2. --- Let \(E_{1}\) and \(E_{2}\) be two reflexive Fréchet spaces, and G a locally convex Hausdorff space. For i = 1, 2, let \(F_{i}\) be the strong dual of \(E_{i}\) . Then every separately continuous bilinear mapping \(u: F_{1} \times F_{2} \to G\) is continuous.

The space G is isomorphic to a subspace of a product of Banach spaces (II, p.~5, prop. 3). Therefore it is enough to prove the theorem under the additional hypothesis that G is a Banach space. But \(F_{1}\) is barrelled and \(F_{2}\) bornological (IV, p.~24, corollary), and \(\mathcal{L}_{b}(F_{2}; G)\) is a Fréchet space (IV, p.~23, prop. 3). Let v denote the linear mapping from \(F_{1}\) into \(\mathcal{L}_{b}(F_{2}, G)\) associated with u by the relation

\[
u (x _ {1}, x _ {2}) = v (x _ {1}) \left(x _ {2}\right) \quad \left(x _ {1} \in \mathrm{F} _ {1}, x _ {2} \in \mathrm{F} _ {2}\right).
\]

Since \(F_{1}\) is barrelled and u separately continuous, v is continuous (III, p.~31, prop. 6). Since v is continuous, lemma 1 implies the existence of a neighbourhood \(U_{1}\) of 0 in \(F_{1}\) whose image under v is bounded in \(\mathcal{L}_{b}(F_{2}; G)\) . In other words, for every bounded subset \(B_{2}\) in \(F_{2}\) , the set \(u(U_{1} \times B_{2})\) is bounded in the Banach space G. Let \(U_{2}\) be the set of all \(x_{2} \in F_{2}\) such that \(\|u(x_{1}, x_{2})\| \leqslant 1\) for all \(x_{1} \in U_{1}\) . The set \(U_{2}\) then absorbs every bounded subset; since \(F_{2}\) is bornological, \(U_{2}\) is a neighbourhood of 0 in \(F_{2}\) , and this proves that u is continuous.

